\section{Near-Collision, Photon Production, and Entanglement}
\paragraph*{Proposition 5.1 (Compensator as near-collision in K3). The Heisenberg compensator C(f) = 1}
2(f +Jf) is the projection onto K3 = e1e2: it takes matter f (in K1) and its tachyon image Jf (in K2), and produces the photon C(f) $\in$K3. The quaternion identity e3 = e1e2 ensures this is always orthogonal to both inputs.
\paragraph*{Theorem 5.2 (Entanglement as shared proper time). Two photons $\phi$1, $\phi$2 $\in$}
K3 are entangled iff they share the same pre-image: $\phi$1 = $\phi$2 = C(f) for some f $\in$K1. The entanglement correlation is the invariant proper time $\tau$n(f) = $\tau$n(Jf) fixed at the moment of the near-collision. No signal is transmitted at measurement: the invariant $\tau$n was established at the near-collision and is merely revealed.
\paragraph*{Proof. $\tau$n(Jf)2 = t2}
n-|-z-|2 = t2 n-|z-|2 = $\tau$n(f)2 by Theorem 2.1. The shared value $\tau$n is a Lorentz invariant at every stratum, hence frame-independent.

